% ENEM 2017-2024 — 9 questões MCQ — Função Afim
% Fonte: lista "Função Afim ou Função do 1º Grau" do site projetoagathaedu.com.br
% Omitidas: Q2,Q3,Q7,Q8,Q11,Q13,Q15,Q18,Q20 (imagens)
%           Q10 (duplicata de lista14/Q8 — soja, L(x)=50x-12000)
%           Q17 (duplicata idêntica de Q16)

---
tipo: Múltipla Escolha
dificuldade: Fácil
nivel: médio
disciplina: Matemática
assunto: Função Afim
gabarito: C
tags: [doceria, docinhos, taxa de entrega, porções]
fonte: concurso
concurso: ENEM
banca: INEP
ano: 2024
numero: 1
---
\question
Uma doceria vende e entrega, em seu bairro, porções de 100 g de docinhos de aniversário. Atualmente, a taxa única de entrega é R\$ 10,00, e o valor cobrado por uma porção é R\$ 25,00. Por uma estratégia de vendas, a partir da próxima semana, a taxa única de entrega será R\$ 15,00, e um novo valor será cobrado por uma porção, de maneira que o valor total a ser pago por um cliente na compra de 5 porções permaneça o mesmo.

A partir da próxima semana, qual será o novo valor cobrado, em real, por uma porção?
\begin{oneparchoices}
\choice 12,50
\choice 20,00
\correctchoice 24,00
\choice 30,00
\choice 37,50
\end{oneparchoices}

---
tipo: Múltipla Escolha
dificuldade: Média
nivel: médio
disciplina: Matemática
assunto: Função Afim
gabarito: D
tags: [piscina, vazão, abastecimento, proporcionalidade]
fonte: concurso
concurso: ENEM
banca: INEP
ano: 2024
numero: 4
---
\question
Uma piscina tem capacidade de 2.500.000 litros. Seu sistema de abastecimento foi regulado para ter uma vazão constante de 6.000 litros de água por minuto.

O mesmo sistema foi instalado em uma segunda piscina, com capacidade de 2.750.000 litros, e regulado para ter uma vazão, também constante, capaz de enchê-la em um tempo 20\% maior que o gasto para encher a primeira piscina.

A vazão do sistema de abastecimento da segunda piscina, em litro por minuto, é
\begin{oneparchoices}
\choice 8.250.
\choice 7.920.
\choice 6.545.
\correctchoice 5.500.
\choice 5.280.
\end{oneparchoices}

---
tipo: Múltipla Escolha
dificuldade: Fácil
nivel: médio
disciplina: Matemática
assunto: Função Afim
gabarito: D
tags: [fotocopiadora, aluguel, cópias, menor custo]
fonte: concurso
concurso: ENEM-PPL
banca: INEP
ano: 2024
numero: 5
---
\question
Uma escola analisou as propostas de cinco empresas para alugar uma máquina fotocopiadora que atenda à demanda de 12.000 cópias mensais. Cada empresa cobra um valor fixo pelo aluguel mensal da máquina, mais um valor proporcional ao número de cópias realizadas, ambos em real. Assim, o custo total \( C \), do aluguel de uma máquina, que atenda a uma demanda de \( x \) cópias mensais, em cada uma das cinco empresas, pode ser dado pelas expressões:

-- empresa I: \( C = 500 + 0{,}40x \);

-- empresa II: \( C = 800 + 0{,}50x \);

-- empresa III: \( C = 2{.}000 + 0{,}20x \);

-- empresa IV: \( C = 1{.}100 + 0{,}25x \);

-- empresa V: \( C = 600 + 0{,}30x \).

A escola escolheu a empresa que apresentou a proposta que fornecia o serviço necessário pelo menor custo mensal.

A empresa escolhida foi a
\begin{oneparchoices}
\choice I.
\choice II.
\choice III.
\correctchoice IV.
\choice V.
\end{oneparchoices}

---
tipo: Múltipla Escolha
dificuldade: Fácil
nivel: médio
disciplina: Matemática
assunto: Função Afim
gabarito: B
tags: [reparos domésticos, prestação de serviço, comparação de custos, inequação]
fonte: concurso
concurso: ENEM-PPL
banca: INEP
ano: 2023
numero: 6
---
\question
Duas empresas do mercado de pequenos reparos domésticos determinam o valor de seus serviços a partir de um valor fixo acrescido de um valor cobrado por hora. A empresa X cobra R\$ 60,00 de valor fixo mais R\$ 18,00 por hora de serviço prestado. A empresa Y cobra R\$ 24,00 de valor fixo e está definindo um novo valor a ser cobrado por hora. Sua estratégia de mercado prevê que, em relação à empresa X, o custo total do serviço deve ser menor ou igual para trabalhos de até duas horas de duração.

Qual é o valor máximo, em real, que a empresa Y poderá cobrar por hora de serviço prestado a fim de atender à sua estratégia de mercado?
\begin{oneparchoices}
\choice 18
\correctchoice 36
\choice 48
\choice 54
\choice 78
\end{oneparchoices}

---
tipo: Múltipla Escolha
dificuldade: Fácil
nivel: médio
disciplina: Matemática
assunto: Função Afim
gabarito: A
tags: [hospedagem, diária, taxa de serviço, aplicativo]
fonte: concurso
concurso: ENEM
banca: INEP
ano: 2021
numero: 9
---
\question
Aplicativos que gerenciam serviços de hospedagem têm ganhado espaço no Brasil e no mundo por oferecer opções diferenciadas em termos de localização e valores de hospedagem. Em um desses aplicativos, o preço \( P \) a ser pago pela hospedagem é calculado considerando um preço por diária \( d \), acrescido de uma taxa fixa de limpeza \( L \) e de uma taxa de serviço. Essa taxa de serviço é um valor percentual \( s \) calculado sobre o valor pago pelo total das diárias.

Nessa situação, o preço a ser pago ao aplicativo para uma hospedagem de \( n \) diárias pode ser obtido pela expressão
\begin{choices}
\correctchoice \( P = d \cdot n + L + d \cdot n \cdot s \)
\choice \( P = d \cdot n + L + d \cdot s \)
\choice \( P = d + L + s \)
\choice \( P = a \cdot d \cdot n \cdot s + L \)
\choice \( P = d \cdot n + L + s \)
\end{choices}

---
tipo: Múltipla Escolha
dificuldade: Média
nivel: médio
disciplina: Matemática
assunto: Função Afim
gabarito: D
tags: [restaurante, couvert, por quilo, taxa fixa, sistema de equações]
fonte: concurso
concurso: ENEM
banca: INEP
ano: 2020
numero: 12
---
\question
Para aumentar a arrecadação de seu restaurante que cobra por quilograma, o proprietário contratou um cantor e passou a cobrar dos clientes um valor fixo de couvert artístico, além do valor da comida. Depois, analisando as planilhas do restaurante, verificou-se em um dia que 30 clientes consumiram um total de 10 kg de comida em um período de 1 hora, sendo que dois desses clientes pagaram R\$ 50,00 e R\$ 34,00 e consumiram 500 g e 300 g, respectivamente.

Qual foi a arrecadação obtida pelo restaurante nesse período de 1 hora, em real?
\begin{oneparchoices}
\choice 800,00.
\choice 810,00.
\choice 820,00.
\correctchoice 1.100,00.
\choice 2.700,00.
\end{oneparchoices}

---
tipo: Múltipla Escolha
dificuldade: Média
nivel: médio
disciplina: Matemática
assunto: Função Afim
gabarito: E
tags: [chaveiros, brinde, ponto de equilíbrio, microempresa]
fonte: concurso
concurso: ENEM
banca: INEP
ano: 2020
numero: 14
---
\question
Uma microempresa especializou-se em produzir um tipo de chaveiro personalizado para brindes. O custo de produção de cada unidade é de R\$ 0,42 e são comercializados em pacotes com 400 chaveiros, que são vendidos por R\$ 280,00. Além disso, essa empresa tem um custo mensal fixo de R\$ 12.800,00 que não depende do número de chaveiros produzidos.

Qual é o número mínimo de pacotes de chaveiros que devem ser vendidos mensalmente para que essa microempresa não tenha prejuízo no mês?
\begin{oneparchoices}
\choice 26
\choice 46
\choice 109
\choice 114
\correctchoice 115
\end{oneparchoices}

---
tipo: Múltipla Escolha
dificuldade: Média
nivel: médio
disciplina: Matemática
assunto: Função Afim
gabarito: D
tags: [gerente, diarista, folha de pagamento, funcionários]
fonte: concurso
concurso: ENEM
banca: INEP
ano: 2019
numero: 16
---
\question
Uma empresa tem diversos funcionários. Um deles é o gerente, que recebe R\$ 1.000,00 por semana. Os outros funcionários são diaristas. Cada um deles trabalha 2 dias por semana, recebendo R\$ 80,00 por dia trabalhado.

Chamando de \( X \) a quantidade total de funcionários da empresa, a quantia \( Y \), em reais, que esta empresa gasta semanalmente para pagar seus funcionários é expressa por
\begin{choices}
\choice \( Y = 80X + 920 \)
\choice \( Y = 80X + 1{.}000 \)
\choice \( Y = 80X + 1{.}080 \)
\correctchoice \( Y = 160X + 840 \)
\choice \( Y = 160X + 1{.}000 \)
\end{choices}

---
tipo: Múltipla Escolha
dificuldade: Fácil
nivel: médio
disciplina: Matemática
assunto: Função Afim
gabarito: A
tags: [reservatório, torneira, vazão, volume de água]
fonte: concurso
concurso: ENEM
banca: INEP
ano: 2017
numero: 19
---
\question
Um reservatório de água com capacidade para 20 mil litros encontra-se com 5 mil litros de água num instante inicial \( t \) igual a zero, em que são abertas duas torneiras. A primeira delas é a única maneira pela qual a água entra no reservatório, e ela despeja 10 L de água por minuto; a segunda é a única maneira de a água sair do reservatório. A razão entre a quantidade de água que entra e a que sai, nessa ordem, é igual a 5/4.

Considere que \( Q(t) \) seja a expressão que indica o volume de água, em litro, contido no reservatório no instante \( t \), dado em minuto, com \( t \) variando de 0 a 7.500.

A expressão algébrica para \( Q(t) \) é
\begin{oneparchoices}
\correctchoice \( 5{.}000 + 2t \)
\choice \( 5{.}000 - 8t \)
\choice \( 5{.}000 - 2t \)
\choice \( 5{.}000 + 10t \)
\choice \( 5{.}000 - 2{,}5t \)
\end{oneparchoices}
